\section{Whether that equation can be believed, and what the code fits instead}

Two different equations are in play, and they must not be confused with each other.

The first is the one asked about. It is ordinary least squares on the nine GO-10 runs, with
$J$ and $R$ in their real units, fitting all five estimable terms at once:

\begin{equation}
  \label{eq:five-term}
  \widehat{\text{recovery}}\;\%
  = 33.0192 + 3.16999\,J - 14.4538\,R - 0.0954576\,J^{2} + 0.837907\,R^{2}.
\end{equation}

The second is the one the code prints. It is a different fit: logged and coded units, and only
the terms that earned their place by forward stepwise selection.

\begin{equation}
  \label{eq:code-fit}
  \widehat{\text{recovery}}\;\%
  = 11.5776 - 15.627\,x_2 + 6.08016\,x_2^{2},
  \qquad x_2 = \frac{\ln R - 1.395778}{1.039721}.
\end{equation}

Equation~\eqref{eq:five-term} is a hand fit, made to answer the question ``what do five terms in
real units come to''. Equation~\eqref{eq:code-fit} is what \texttt{uv run limodel analyze} prints
today. They share the nine runs and nothing else: different units, different terms, different
column space. Current density does not appear in~\eqref{eq:code-fit} at all. The rest of this
section judges the first and then explains the second.

\subsection{$R^{2}$ says 0.875, and $R^{2}$ cannot say otherwise}

$R^{2}$ compares the squared misses of the fit against the squared misses of the flat line at
the mean. For~\eqref{eq:five-term} the sum of squared residuals is $201.9127$ and the same sum
about the mean of the nine recoveries is $1611.1333$, so

\begin{equation}
  R^{2} = 1 - \frac{201.9127}{1611.1333} = 0.875.
\end{equation}

That number is not evidence, because of how least squares works. Adding a column to $X$ widens
the space the fit may land in, and it still contains everything the narrower space contained ---
setting the new coefficient to zero recovers the old fit exactly. Least squares takes the best
point in the wider space, so the sum of squared residuals can only fall and $R^{2}$ can only
rise. A term that is pure noise still buys a little. With nine runs and five columns there is
plenty to buy.

Table~\ref{tab:r2-ladder} adds the real-unit terms one at a time and prints $R^{2}$ beside the
leave-one-out error. $R^{2}$ climbs every single step. The out-of-sample error does not.

\begin{table}[htbp]
  \centering
  \caption{The nested real-unit ladder on the nine GO-10 runs. $R^{2}$ rises by construction;
  the leave-one-out error is free to do what it likes, and gets worse.}
  \label{tab:r2-ladder}
  \begin{tabular}{@{}lrrr@{}}
    \toprule
    Terms in the fit & Residual sum of squares & $R^{2}$ & Leave-one-out rmse \\
    \midrule
    intercept only              & 1611.1333 & 0.000 & 15.05 \\
    ${}+J$                      & 1598.2892 & 0.008 & 16.31 \\
    ${}+R$                      &  507.1461 & 0.685 & 15.56 \\
    ${}+J^{2}$                  &  505.2442 & 0.686 & 19.82 \\
    ${}+R^{2}$                  &  201.9127 & 0.875 & 41.77 \\
    \bottomrule
  \end{tabular}
\end{table}

\subsection{Leave-one-out: fit on eight runs, predict the ninth}

The honest question is not how close the curve passes to the runs it was fitted on, but how far
it would miss a run it has never seen. With nine runs there is no test set to spare, so every run
takes its turn as the test set: drop run~$i$, fit the same five terms on the remaining eight,
predict run~$i$, square the miss. Nine folds, then the root of the mean.

For~\eqref{eq:five-term} that comes to

\begin{equation}
  \text{rmse}_{\text{LOO}} = 41.77 \;\text{points of recovery}.
\end{equation}

Three yardsticks sit beside it, and all three are smaller:

\begin{itemize}[leftmargin=1.4em]
  \item \textbf{6.612} --- the bench's own scatter, the pooled standard deviation of the two
    settings that were run twice. This is the floor. No equation can predict a run more closely
    than the bench reproduces one.
  \item \textbf{14.19} --- the standard deviation of the nine recoveries. This is what you miss
    by if you ignore both settings and quote the average, $25.14\,\%$, for every run.
  \item \textbf{15.05} --- the leave-one-out error of quoting that average, which is the same
    idea measured the same way as the 41.77, fold for fold.
\end{itemize}

So the conclusion is not ``the fit is weak''. It is worse than that. Predicting a new run with
\eqref{eq:five-term} misses by $41.77$ where quoting the average of the nine misses by $15.05$,
a factor of $2.77$. The equation has \emph{negative skill}: on the evidence available it is
strictly worse than refusing to model at all.

Where that comes from is countable. Five coefficients are being fitted to nine runs, which
leaves four degrees of freedom. Worse, the runs are not spread evenly across the two terms in
$R$: six of the nine (AE, AH, AI, AJ, AN, AP) sit at the lowest feed ratio, $1.4277$, and the
feed ratio moves in only three runs --- AF, AG and AT. Those three runs are the sole support for
both the $R$ term and the $R^{2}$ term. Drop any one of them in a fold and the curvature in $R$
is being guessed from two points, which is why the leave-one-out error is so much larger than
the fitted residuals.

\subsection{The two repeats, and what they cap}

Two settings were run twice. The code finds them in
\texttt{ReplicateScatter\allowbreak{}.find\_groups}, which groups runs by membrane and both
settings rounded to two places.

\begin{table}[htbp]
  \centering
  \caption{The repeated settings on GO-10. Both pairs are the same dialled setup, run twice.}
  \label{tab:repeats}
  \begin{tabular}{@{}lrrrrr@{}}
    \toprule
    Pair & $J$ ($\mathrm{mA/cm^2}$) & $R$ ($\mathrm{mol/mol}$) & First & Second & Difference \\
    \midrule
    AH, AI & 7.50  & 1.4277 & 36.59 & 37.66 & 1.07 \\
    AN, AP & 25.67 & 1.4277 & 39.50 & 26.32 & 13.18 \\
    \bottomrule
  \end{tabular}
\end{table}

Pooling the two pairs gives the $6.612$ quoted above, on two degrees of freedom. The second pair
is doing almost all of the work: AN and AP are the same setting and disagree by $13.2$ points of
recovery, which is nearly as wide as the whole spread of the nine runs.

This caps every fit, not just this one. A curve cannot resolve an effect smaller than the
disagreement between two runs of the same thing. Any claim that moving a setting changes recovery
by, say, eight points is indistinguishable from running the same setting twice. It also distorts
the fit: in~\eqref{eq:five-term} AN and AP get the same prediction, $32.57$, and so contribute
residuals of $+6.93$ and $-6.25$ no matter what the coefficients are. Those two alone account
for $87$ of the $201.9$ residual sum of squares.

\begin{table}[htbp]
  \centering
  \caption{Fitted values and residuals of~\eqref{eq:five-term}. Residual is measured recovery
  minus fitted.}
  \label{tab:five-term-residuals}
  \begin{tabular}{@{}lrrrr@{}}
    \toprule
    Run & $J$ ($\mathrm{mA/cm^2}$) & $R$ ($\mathrm{mol/mol}$) & Fitted \% & Residual \\
    \midrule
    AE & 13.67 &  1.4277 & 39.59 & $-1.31$ \\
    AF & 13.67 &  2.8554 & 24.07 & $-5.73$ \\
    AG & 13.67 &  5.7108 &  3.30 & $+2.86$ \\
    AH &  7.50 &  1.4277 & 32.50 & $+4.09$ \\
    AI &  7.50 &  1.4277 & 32.50 & $+5.16$ \\
    AJ &  4.50 &  1.4277 & 26.42 & $-5.34$ \\
    AN & 25.67 &  1.4277 & 32.57 & $+6.93$ \\
    AP & 25.67 &  1.4277 & 32.57 & $-6.25$ \\
    AT & 13.67 & 11.4215 &  2.73 & $-0.41$ \\
    \midrule
    \multicolumn{4}{@{}l}{Residual sum of squares} & 201.9127 \\
    \bottomrule
  \end{tabular}
\end{table}

\subsection{Why there is no interaction term}

The obvious sixth term is $J\!\cdot\!R$, the one that would say whether the best current density
depends on how much magnesium is in the feed. It cannot be fitted. Not fitted badly --- not
fitted at all.

Over these nine runs, exactly,

\begin{equation}
  \label{eq:no-interaction-term}
  J\!\cdot\!R = 1.4277\,J + 13.6667\,R - 19.5118,
\end{equation}

and ``exactly'' is meant literally: the largest departure over the nine rows is
$7.1\times10^{-15}$, which is floating-point dust. Put the six columns side by side,

\begin{equation}
  X = \begin{bmatrix}
    1 & J_1 & R_1 & J_1^2 & R_1^2 & J_1 R_1 \\
    \vdots & \vdots & \vdots & \vdots & \vdots & \vdots \\
    1 & J_9 & R_9 & J_9^2 & R_9^2 & J_9 R_9
  \end{bmatrix},
\end{equation}

and the rank is $5$, not $6$. The interaction column is already a combination of the intercept,
the $J$ column and the $R$ column. There is no unique answer for the six coefficients: shift the
interaction coefficient by $t$, adjust the intercept by $19.5118\,t$, the $J$ coefficient by
$-1.4277\,t$ and the $R$ coefficient by $-13.6667\,t$, and you have a different equation that
fits every one of the nine runs identically well. This is why
\texttt{SecondOrderPolynomial\allowbreak{}.find\_fit} checks the rank \texttt{numpy.linalg.lstsq}
returns and hands back \texttt{None} when it falls short of the number of terms, instead of
printing whichever of the infinitely many answers the solver happened to land on.

The reason is the shape of the design, not the arithmetic. Figure~\ref{fig:design-cross} draws
it. Current density was moved only at the lowest feed ratio, and the feed ratio was moved only at
the middle current density. The runs form a cross, with AE at the crossing point. On a cross,
one of the two settings is always sitting at its base value, so for every run
$(J - 13.6667)(R - 1.4277) = 0$. Expanding that product gives
\eqref{eq:no-interaction-term} directly, with $13.6667 \times 1.4277 = 19.5118$. The product
column carries nothing the other columns do not already carry. Fitting an interaction needs at
least one run off the cross --- a high feed ratio at a high current density, say.

\begin{figure}[htbp]
  \centering
  \begin{tikzpicture}[x=4.6cm, y=1.35cm, font=\small]
    % axes; points are placed on log spacing, which is how the code sees them
    \draw[-Stealth] (-1.25,-1.3) -- (1.3,-1.3);
    \node at (0.1,-1.95) {$J$ ($\mathrm{mA/cm^2}$), log spacing};
    \draw[-Stealth] (-1.25,-1.3) -- (-1.25,1.35) node[above, align=center]
      {$R$ ($\mathrm{mol/mol}$),\\log spacing};

    % tick labels
    \foreach \x/\lab in {-1/4.5, -0.4132/7.5, 0.2761/13.67, 1/25.67}
      \draw (\x,-1.3) -- (\x,-1.38) node[below, font=\footnotesize] {\lab};
    \foreach \y/\lab in {-1/1.43, -0.3333/2.86, 0.3333/5.71, 1/11.42}
      \draw (-1.25,\y) -- (-1.31,\y) node[left, font=\footnotesize] {\lab};

    % the arm in current density, all at the lowest feed ratio
    \draw[gray, dashed] (-1,-1) -- (1,-1);
    % the arm in feed ratio, all at the middle current density
    \draw[gray, dashed] (0.2761,-1) -- (0.2761,1);

    \foreach \x/\y/\lab in {%
      -1/-1/AJ, -0.4132/-1/{AH, AI}, 1/-1/{AN, AP},
      0.2761/-0.3333/AF, 0.2761/0.3333/AG, 0.2761/1/AT}
      \fill (\x,\y) circle (2.2pt) node[above right=1pt and 2pt, font=\footnotesize] {\lab};
    \fill (0.2761,-1) circle (2.2pt)
      node[above right=1pt and 2pt, font=\footnotesize] {AE};
  \end{tikzpicture}
  \caption{The nine GO-10 runs. Current density moves only along the bottom arm, at
  $R = 1.4277$; the feed ratio moves only along the upright arm, at $J = 13.67$. AE sits where
  they cross, and AH/AI and AN/AP are the two repeated settings. Nothing occupies the interior,
  so the $J\!\cdot\!R$ column is redundant.}
  \label{fig:design-cross}
\end{figure}

\subsection{What the code fits instead}

The code never forms~\eqref{eq:five-term}. It makes two choices the hand fit does not.

\paragraph{Logged, then coded.} \texttt{DesignPoints} takes the logarithm of each setting, then
centres and scales it so the runs span $-1$ to $+1$:

\begin{align}
  x_1 &= \frac{\ln J - 2.374635}{0.870558},
  & c_J &= \tfrac12(\ln J_{\max} + \ln J_{\min}),
  & h_J &= \tfrac12(\ln J_{\max} - \ln J_{\min}), \\
  x_2 &= \frac{\ln R - 1.395778}{1.039721},
  & c_R &= \tfrac12(\ln R_{\max} + \ln R_{\min}),
  & h_R &= \tfrac12(\ln R_{\max} - \ln R_{\min}).
\end{align}

Logs because the settings were dialled multiplicatively, not in equal steps: $J$ runs
$4.5 \to 25.67$, a $5.7$-fold span, and $R$ runs $1.4277 \to 11.42$, which is three exact
doublings, an $8$-fold span. A polynomial in $R$ asked to cover an eightfold range spends its
curvature on the top end.

Coding because without it the squared columns swamp the linear ones. In real units the $J$
column reaches $25.67$ while $J^{2}$ reaches $658.8$, and $R$ reaches $11.42$ while $R^{2}$
reaches $130.5$. The condition number of the five-column $X$ is $1770$. Logged but uncoded it is
$225$. Logged and coded, every column lies within $\pm 1$ and the condition number is $6.07$.
That is not cosmetic: a badly scaled $X$ makes the coefficients delicate, and coefficients that
move a lot for a small change in the data are coefficients nobody should quote.

\paragraph{Forward stepwise on the leave-one-out error.} Rather than fit everything and report
it, \texttt{FittedOutcome\allowbreak{}.find\_best\_polynomial} grows the model one term at a
time, and the thing it minimises is the out-of-sample error, not the fitted residual. That
matters, because minimising the fitted residual would add every term every time, as
Table~\ref{tab:r2-ladder} shows. Figure~\ref{fig:stepwise} is the loop.

The one structural rule is in \texttt{find\_terms\_it\_may\_add}: a squared term may not enter
before its own linear term. A curvature without its slope cannot be read --- the vertex is then
pinned to the centre of the coded range by accident rather than by evidence.

\begin{figure}[htbp]
  \centering
  \begin{tikzpicture}[font=\small, node distance=0pt,
      box/.style={draw, rounded corners, align=center, inner sep=5pt, text width=8.2cm},
      arrow/.style={-Stealth, thick}]

    \node[box] (start) at (0,0)
      {\textbf{Start.} chosen $=$ (intercept).\\
       $e \leftarrow$ leave-one-out rmse of the intercept-only fit};

    \node[box, below=0.65cm of start] (allowed)
      {List the terms it may add: those from
       $(x_1,\; x_2,\; x_1^{2},\; x_2^{2})$ not yet chosen, whose own linear term is already in};

    \node[box, below=0.65cm of allowed] (score)
      {For each one, fit chosen $+$ that term nine times over, leaving out one run each time,
       and take the rmse of the nine misses\\[2pt]
       \footnotesize (a fit whose rank falls short of its terms returns no rmse; that term is
       passed over)};

    \node[box, below=0.65cm of score] (best)
      {Take the lowest rmse of them, $r$, and the term that reached it};

    \node[draw, rounded corners, below=0.65cm of best, inner sep=5pt,
          minimum width=8.4cm] (test)
      {Is $r < e$?};

    \node[box, below=1.15cm of test] (stop)
      {\textbf{Stop.} Refit chosen on all nine runs and print those coefficients.\\
       Report $e$ as the fit's leave-one-out error};

    \draw[arrow] (start) -- (allowed);
    \draw[arrow] (allowed) -- (score);
    \draw[arrow] (score) -- (best);
    \draw[arrow] (best) -- (test);
    \draw[arrow] (test) -- node[right, font=\footnotesize]
      {no, or nothing left to add} (stop);

    % the loop back up the left-hand side
    \coordinate (left) at ($(test.west) + (-1.1cm, 0)$);
    \draw[arrow] (test.west) -- (left) -- (left |- allowed.west) -- (allowed.west);
    \node[font=\footnotesize, align=center, left=1pt of left]
      {yes:\\ add it,\\ $e \leftarrow r$};
  \end{tikzpicture}
  \caption{Forward stepwise selection as \texttt{FittedOutcome.find\_best\_polynomial} runs it.
  The loop adds a term only when it lowers the leave-one-out error, so a term that merely fits
  the nine runs better never gets in.}
  \label{fig:stepwise}
\end{figure}

Run on recovery, the loop takes two steps and stops (Table~\ref{tab:stepwise-trace}). Current
density is offered at every step and rejected at every step: adding $x_1$ always makes the
out-of-sample error worse. Given that AN and AP disagree by $13.2$ points at the same current
density, that is the honest outcome.

\begin{table}[htbp]
  \centering
  \caption{The stepwise trace on recovery \%. Each cell is a full nine-fold leave-one-out rmse.
  Candidates are offered in the order $x_1$, $x_2$, $x_1^2$, $x_2^2$, so a tie goes to the
  simpler term.}
  \label{tab:stepwise-trace}
  \begin{tabular}{@{}llrl@{}}
    \toprule
    Step & Model so far & Its rmse & Terms offered, and their rmse \\
    \midrule
    0 & intercept                     & 15.052 & $x_1$: 16.244 \quad $x_2$: \textbf{7.898} \\
    1 & intercept, $x_2$              &  7.898 & $x_1$: 9.950 \quad $x_2^2$: \textbf{7.122} \\
    2 & intercept, $x_2$, $x_2^{2}$   &  7.122 & $x_1$: 9.471 --- nothing lower, stop \\
    \bottomrule
  \end{tabular}
\end{table}

That leaves~\eqref{eq:code-fit}, which the report prints as
\texttt{intercept +11.58,\allowbreak{} x2 -15.63,\allowbreak{} x2\_squared +6.08}. Its leave-one-out error is $7.122$
against the same bench scatter of $6.612$.

This is a far better equation than~\eqref{eq:five-term}: $7.122$ against $41.77$, and now
comfortably inside the $15.05$ of quoting the average, so it does carry real information about
the feed ratio. But $7.122$ is still larger than $6.612$, and that is the comparison that
decides. The fit misses a held-out run by more than the bench misses itself when it repeats a
setting, so nothing in the equation has been separated from the bench's own noise. The report
says so in those words --- \emph{worse than the bench repeating itself, so this fit has earned
no conclusion} --- and for the same reason the search under the gates returns no best setup: no
point inside the measured range clears the $35\,\%$ recovery floor.

\subsection{What would change the answer}

Two things would, and they are both runs rather than arithmetic.

Resolving AN against AP settles the yardstick. Those two are the same setting at
$25.67\unit{mA/cm^2}$ and $1.4277\unit{mol/mol}$ and they differ by $13.2$ points of recovery.
That single disagreement sets the $6.612$ that every fit is measured against, and until it is
either reproduced or explained, no fit can clear it. A third run at that setting either confirms
the spread, in which case recovery at high current density is genuinely that variable, or
identifies one of the two as faulty, in which case the scatter drops and $7.122$ may already be
good enough.

Runs off the cross make the interaction estimable and the curvature in $J$ real. A feed ratio
away from the lowest, at a current density away from the middle --- $R = 5.71$ at
$J = 25.67\unit{mA/cm^2}$, say --- breaks the identity
$(J - 13.6667)(R - 1.4277) = 0$, lifts the rank of the six-column design from $5$ to $6$, and
moves the support for the terms in $R$ off the three runs AF, AG and AT that currently carry
them alone.
